% !PRINTED 170
\weberpage{170}

But we have
\[
F(\tfrac12,\tfrac12,1\pm\varepsilon,x) = \frac{\Pi(\pm\varepsilon)}{[\Pi(-\tfrac12)]^2} \sum_{\nu=0}^{\infty} \frac{\Pi(\nu-\tfrac12)^2}{\Pi(\nu\pm\varepsilon)\Pi(\nu)} x^\nu,
\]
and, with the addition of a constant factor, one may also take
the following as the second particular integral:
\[
\frac{1}{2\varepsilon}\cdot\sum_{\nu=0}^{\infty} \frac{\Pi(\nu-\tfrac12)^2}{\Pi(\nu)}\Bigl[\frac{-1}{\Pi(\nu+\varepsilon)} + \Bigl(\frac{x}{1-x}\Bigr)^{-\varepsilon}\frac{1}{\Pi(\nu-\varepsilon)}\Bigr]x^\nu.
\]
If we here develop in powers of $\varepsilon$ and then set
$\varepsilon=0$, it follows that:
\[
\sum_{\nu=0}^{\infty} \frac{\Pi(\nu-\tfrac12)^2}{\Pi(\nu)^2} x^\nu \Bigl[\frac{\Pi'(\nu)}{\Pi(\nu)} - \frac12 \log\frac{x}{1-x}\Bigr],
\]
for which, using the \gesperrt{Gauss} formula
\[
\frac{\Pi(\nu-\tfrac12)\Pi(\nu)}{\Pi(2\nu)} = \sqrt{\pi}\,2^{-2\nu}
\]
one may also write, on discarding a constant factor:
\weq{12}{\sum_{\nu=0}^{\infty} \frac{\Pi(2\nu)^2}{\Pi(\nu)^4}\Bigl(\frac{x}{16}\Bigr)^\nu \Bigl[\frac{\Pi'(\nu)}{\Pi(\nu)} - \frac12 \log\frac{x}{1-x}\Bigr].}
But for a positive integer $\nu$ we have:
\[
\frac{\Pi'(\nu)}{\Pi(\nu)} = \Pi'(0) + 1 + \frac12 + \frac13 + \cdots + \frac1\nu,
\]
and if we therefore set, for brevity:
% [sic] In the second line the source omits the plus sign between \cdots and 1/\nu.
\weq{13}{\begin{aligned}
F(x) &= \sum_{\nu=0}^{\infty} \frac{\Pi(2\nu)^2}{\Pi(\nu)^4}\Bigl(\frac{x}{16}\Bigr)^\nu, \\
G(x) &= 2\sum_{\nu=1}^{\infty} \frac{\Pi(2\nu)^2}{\Pi(\nu)^4}\Bigl(1+\frac12+\frac13+\cdots\frac1\nu\Bigr)\Bigl(\frac{x}{16}\Bigr)^\nu,
\end{aligned}}
then the two particular solutions of the differential equation
 (8) result, if for the second one takes a linear combination
of (10) and (12):
\weq{14}{y_1 = F(x), \qquad y_2 = \frac12 G(x) - \frac12 \log\frac{x}{16(1-x)}F(x).}
But in order now to represent the particular integrals $K$ and $K'$ by
these functions, we must investigate the behaviour of $K,K'$
for $q=0$, for which at the same time $\varkappa^2=0$.

But for $q=0$ we have, by \S\,25 and \S\,42:
\[
q^{-1}\varkappa^2 = 16, \qquad 2K = \pi,
\]
%CONT>>
